% capitulo8_conclusiones.tex
% Capitulo 8: Conclusiones y trabajo futuro
%
% Autor: Leon Elliott Fuller
% Fecha: Mayo 2026

\chapter{Conclusiones y trabajo futuro}
\label{ch:conclusiones}

\section{Conclusiones generales}
\label{sec:conc-generales}

El objetivo central de este Trabajo de Fin de Grado era diseñar, implementar y evaluar de forma reproducible una comparativa de rendimiento entre la criptografía RSA y la criptografía de curvas elípticas (ECC), articulada a lo largo de tres ejes ortogonales: RSA frente a ECC a niveles de seguridad equivalentes, coordenadas afines frente a Jacobianas, y cuerpos primos $\mathbb{F}_p$ frente a cuerpos binarios $\mathbb{F}_{2^m}$. A la luz de los resultados presentados en el capítulo~\ref{ch:resultados}, puede afirmarse que este objetivo se ha alcanzado.

El trabajo combina dos componentes que se refuerzan mutuamente. Por un lado, un desarrollo teórico que recorre los fundamentos de los cuerpos finitos (capítulo~\ref{chap:Campos_finitos}), la teoría de curvas elípticas (capítulo~\ref{chap:Curvas_elipticas}) y los esquemas criptográficos que se construyen sobre ellas, junto con los ataques que determinan su seguridad (capítulo~\ref{ch:ecc}). Por otro, una implementación completa en C++17 (capítulo~\ref{ch:implementacion}) que materializa esa teoría en una suite de benchmarking funcional, capaz de medir RSA, ECC sobre cuerpos primos en dos sistemas de coordenadas, y ECC sobre cuerpos binarios, todo ello sobre una única base aritmética (NTL/GMP) y en un entorno reproducible mediante Docker.

\section{Principales hallazgos}
\label{sec:conc-hallazgos}

Los experimentos permiten extraer varias conclusiones concretas y cuantificadas:

\paragraph{La superioridad de ECC sobre RSA es matizada, no absoluta.}
El resultado más relevante del trabajo es que la idea extendida de que ``ECC siempre es más rápida que RSA'' resulta ser una simplificación. Los datos muestran que ECC domina con claridad en la generación de claves (en torno a 150 veces más rápida al mismo nivel de seguridad) y en la firma (unas 2,6 veces), pero RSA conserva una ventaja sustancial en la verificación (cerca de 38 veces más rápida), gracias a su exponente público pequeño. La conclusión práctica es que la elección entre ambos sistemas debe guiarse por el perfil de uso de la aplicación: ECC es preferible cuando predominan la generación de claves y la firma, mientras que RSA puede seguir siendo competitivo cuando predomina la verificación. No obstante, la enorme ventaja de ECC en tamaño de clave y generación, unida a su mejor escalado al aumentar la seguridad, explica su adopción creciente como estándar de próxima generación.

\paragraph{Las coordenadas Jacobianas aportan una mejora real pero acotada.}
La transición de coordenadas afines a Jacobianas, motivada por la eliminación de la inversión modular en cada operación de grupo, produce una aceleración consistente de entre $1{,}7\times$ y $1{,}9\times$, mayor cuanto más grande es el cuerpo. Este valor es más modesto que el que sugieren los modelos teóricos basados en el conteo de operaciones, y la discrepancia constituye en sí misma un hallazgo valioso: pone de manifiesto que el coste real de la inversión modular en una biblioteca optimizada como NTL/GMP es inferior al supuesto en dichos modelos, de modo que las predicciones teóricas deben matizarse con mediciones empíricas sobre la plataforma concreta.

\paragraph{La ventaja de los cuerpos primos en software es un artefacto de implementación.}
En el entorno de software puro empleado, los cuerpos primos resultan aproximadamente $1{,}8$ veces más rápidos que los binarios al mismo nivel de seguridad. Sin embargo, este resultado no refleja una superioridad intrínseca de $\mathbb{F}_p$, sino el distinto grado de optimización de la aritmética subyacente en NTL (\texttt{ZZ\_p} sobre GMP frente a \texttt{GF2E}). En plataformas con soporte hardware para la multiplicación sin acarreo, la relación podría invertirse. La lección es que las comparaciones de rendimiento entre tipos de cuerpo deben interpretarse siempre en el contexto de la plataforma de ejecución.

\paragraph{Una implementación educativa puede ser competitiva.}
La comparación con OpenSSL revela que la implementación de RSA del proyecto iguala en prestaciones a una biblioteca de producción, por compartir el motor aritmético GMP. En ECC, la brecha con OpenSSL es grande para P-256 (por el ensamblador específico que esta incorpora) pero se reduce notablemente para P-384, lo que demuestra que la diferencia no reside en el algoritmo sino en el ajuste a bajo nivel. Esto confirma que el código desarrollado constituye una base de referencia razonable para fines comparativos y didácticos.

\section{Lecciones aprendidas}
\label{sec:conc-lecciones}

Más allá de los resultados numéricos, el desarrollo del trabajo ha dejado varias enseñanzas de carácter metodológico:

\begin{description}
    \item[La progresión pedagógica como método.] Implementar primero la aritmética en coordenadas afines y solo después añadir las Jacobianas como optimización explícita resultó ser una decisión acertada. Disponer de una implementación de referencia, sencilla y verificable, permitió validar la corrección antes de optimizar, y convirtió la propia optimización en un resultado experimental medible.

    \item[El valor de implementar desde cero.] Desarrollar SHA-256 siguiendo el estándar FIPS~PUB~180-4, en lugar de delegar en una biblioteca externa, no solo aportó comprensión profunda del algoritmo, sino que eliminó una variable externa de las mediciones, contribuyendo a la reproducibilidad.

    \item[La reproducibilidad como requisito, no como añadido.] El uso de semillas fijas en el generador de números aleatorios y la contenedorización con Docker fueron decisiones tomadas desde el inicio, y resultaron esenciales para obtener mediciones estables y comparables. La baja varianza observada en la mayoría de las operaciones es consecuencia directa de este enfoque.

    \item[La teoría se valida con la práctica.] Varias predicciones teóricas se confirmaron experimentalmente (la verificación ECDSA cuesta el doble que la firma, la eficacia del CRT en RSA, el escalado con el tamaño del cuerpo), pero otras requirieron matización a la luz de los datos (el \textit{speedup} Jacobiano, la comparación entre tipos de cuerpo). Esta interacción entre modelo y medición es uno de los aprendizajes más valiosos del trabajo.
\end{description}

\section{Limitaciones del trabajo}
\label{sec:conc-limitaciones}

Conviene reconocer con honestidad las limitaciones del trabajo, muchas de las cuales ya se detallaron en la sección~\ref{sec:impl-limitaciones}:

\begin{enumerate}
    \item \textbf{Carácter educativo, no de producción.} La implementación es un prototipo de benchmarking y no un sistema criptográfico seguro para uso real. Carece de protecciones frente a ataques de canal lateral, emplea RSA sin relleno (\textit{textbook RSA}), genera nonces ECDSA de forma aleatoria en lugar de determinista (RFC~6979), y no aplica un KDF estándar a los secretos ECDH.

    \item \textbf{Alcance de las mediciones.} Los benchmarks se ejecutaron con 20 iteraciones por operación sobre una única máquina. Aunque suficiente para apreciar las tendencias, un número mayor de iteraciones y la ejecución en varias plataformas reforzarían la robustez estadística, especialmente para la generación de claves de RSA, de varianza elevada.

    \item \textbf{Cobertura de la comparación.} El nivel de 192~bits no pudo compararse de forma directa entre RSA y ECC, al resultar prohibitiva la generación de claves RSA de 7680~bits. Asimismo, las curvas binarias solo se evaluaron en coordenadas afines, sin una variante proyectiva análoga a la Jacobiana de los cuerpos primos.

    \item \textbf{Dependencia de la biblioteca aritmética.} Varios resultados (el \textit{speedup} Jacobiano, la comparación entre cuerpos) están condicionados por las características concretas de NTL y GMP, y podrían diferir con otras bibliotecas o con soporte hardware específico.
\end{enumerate}

\section{Trabajo futuro}
\label{sec:conc-futuro}

El trabajo abre varias líneas de continuación naturales, que pueden agruparse en tres categorías:

\paragraph{Optimización de la aritmética de curva.}
La línea más inmediata consiste en implementar las técnicas de optimización de la multiplicación escalar que quedaron fuera del alcance de este trabajo: la \emph{forma no adyacente} (NAF) y sus variantes con ventana ($w$-NAF), que reducen el número de sumas de puntos; las tablas de precomputación para puntos fijos; y la escalera de Montgomery, que además aporta resistencia a canales laterales. Para las curvas de Koblitz sobre cuerpos binarios, el uso del endomorfismo de Frobenius permitiría sustituir duplicaciones por operaciones de cuadrado, mucho más baratas. Sería igualmente interesante medir el impacto de la instrucción \texttt{PCLMULQDQ} en la aritmética de cuerpos binarios, para contrastar la hipótesis de que la desventaja observada es atribuible a la implementación software.

\paragraph{Robustez criptográfica.}
Una segunda línea consistiría en elevar la implementación desde un prototipo educativo hacia un sistema más realista, incorporando los esquemas de relleno OAEP (para cifrado RSA) y PSS (para firma RSA), los nonces deterministas de RFC~6979 para ECDSA, y funciones de derivación de clave estándar como HKDF para los secretos ECDH. Añadir curvas modernas como Curve25519 y Ed25519, diseñadas específicamente para ser rápidas y resistentes a errores de implementación, ampliaría el alcance de la comparación con los estándares actuales.

\paragraph{Ampliación de la comparativa.}
Finalmente, el marco de benchmarking desarrollado podría extenderse para incorporar la dimensión post-cuántica, comparando RSA y ECC con los algoritmos recientemente estandarizados de criptografía resistente a ordenadores cuánticos (como los basados en retículos). Dado que tanto RSA como ECC son vulnerables al algoritmo de Shor, esta comparación resulta especialmente pertinente para valorar el coste de la transición hacia la criptografía post-cuántica, y conectaría de forma natural con la motivación expuesta en la introducción de este trabajo.

\section{Reflexión final}
\label{sec:conc-reflexion}

Este trabajo ha recorrido el camino que va desde los fundamentos algebraicos de los cuerpos finitos hasta la medición empírica del rendimiento de los criptosistemas que se construyen sobre ellos. El resultado es doble: una memoria que documenta la teoría de manera autocontenida y una herramienta de software que permite explorar, de forma cuantificable y reproducible, las consecuencias prácticas de esa teoría. La principal aportación no es ningún dato aislado, sino la comprensión integrada de que las decisiones criptográficas ---qué sistema, qué coordenadas, qué cuerpo--- son siempre un compromiso entre seguridad, rendimiento y contexto de ejecución, y que solo la combinación de análisis teórico y medición rigurosa permite tomarlas con fundamento.